\documentclass{ximera}

\input{../../../preamble.tex}


\definecolor{fvdnavy}{HTML}{071D43}
\definecolor{fvdcyan}{HTML}{20BFC6}
\definecolor{fvdcyanlight}{HTML}{EFFCFB}
\definecolor{fvdmagenta}{HTML}{F10D69}
\definecolor{fvdmagentalight}{HTML}{FFF1F7}
\definecolor{fvdtext}{HTML}{163044}





%=================================================
% Pedagogische omgevingen
% PDF: tcolorbox
% HTML: learning-box
%=================================================

\newenvironment{voorkennis}
{%
    \ifdefined\HCode
        \HCode{%
<section class="learning-box learning-box--prior">
<div class="learning-box__header">
<span class="learning-box__icon" aria-hidden="true"></span>
<span class="learning-box__title">Voorkennis</span>
</div>
<div class="learning-box__content">
        }%
        \begin{itemize}
    \else
       \begin{tcolorbox}[
    enhanced,
    breakable,
    colback=fvdcyanlight,
    colframe=fvdcyan,
    coltext=fvdtext,
    boxrule=0.5pt,
    leftrule=4pt,
    arc=4mm,
    outer arc=4mm,
    left=7mm,
    right=6mm,
    top=5mm,
    bottom=5mm,
    before skip=6mm,
    after skip=6mm,
    title=Voorkennis,
    coltitle=fvdcyan,
    fonttitle=\bfseries\Large,
    attach title to upper={\par\smallskip},
    boxed title style={
        empty,
        left=0mm,
        right=0mm,
        top=0mm,
        bottom=0mm
    },
    overlay={
        \draw[
            fvdcyan,
            line width=0.5pt
        ]
        ([xshift=42mm,yshift=-13mm]frame.north west)
        --
        ([xshift=-7mm,yshift=-13mm]frame.north east);
    }
]
        \begin{itemize}
    \fi
}
{%
    \end{itemize}
    \ifdefined\HCode
        \HCode{%
</div>
</section>
        }%
    \else
        \end{tcolorbox}
    \fi
}


\newenvironment{leerdoelen}
{%
    \ifdefined\HCode
        \HCode{%
<section class="learning-box learning-box--goals">
<div class="learning-box__header">
<span class="learning-box__icon" aria-hidden="true"></span>
<span class="learning-box__title">Leerdoelen</span>
</div>
<div class="learning-box__content">
        }%
        \begin{itemize}
    \else
\begin{tcolorbox}[
    enhanced,
    breakable,
    colback=fvdmagentalight,
    colframe=fvdmagenta,
    coltext=fvdtext,
    boxrule=0.5pt,
    leftrule=4pt,
    arc=4mm,
    outer arc=4mm,
    left=7mm,
    right=6mm,
    top=5mm,
    bottom=5mm,
    before skip=6mm,
    after skip=6mm,
    title=Leerdoelen,
    coltitle=fvdmagenta,
    fonttitle=\bfseries\Large,
    attach title to upper={\par\smallskip},
    boxed title style={
        empty,
        left=0mm,
        right=0mm,
        top=0mm,
        bottom=0mm
    },
    overlay={
        \draw[
            fvdmagenta,
            line width=0.5pt
        ]
        ([xshift=42mm,yshift=-13mm]frame.north west)
        --
        ([xshift=-7mm,yshift=-13mm]frame.north east);
    }
]
        \begin{itemize}
    \fi
}
{%
    \end{itemize}
    \ifdefined\HCode
        \HCode{%
</div>
</section>
        }%
    \else
        \end{tcolorbox}
    \fi
}


\addPrintStyle{../../..}

\title{De kettingregel}
\author{Ingmar Herreman}

\begin{abstract}
Niet elke functie bestaat uit een eenvoudige som, een product of een quotiënt.
Vaak wordt een functie in een andere functie ingevuld.

Zo is
\[(3x-1)^5\]
geen gewone macht van $x$: eerst wordt $3x-1$ berekend en daarna wordt
het resultaat tot de vijfde macht verheven.

Zulke functies noemen we samengestelde functies.

In dit hoofdstuk leer je hun structuur herkennen en afleiden met de
kettingregel. Daarbij staat niet alleen het uitvoeren van de regel centraal,
maar vooral het herkennen van de verschillende functielagen en het kiezen
van een efficiënte afleidingsstrategie.

De kettingregel wordt later opnieuw gebruikt bij rationale, irrationale,
exponentiële, logaritmische, goniometrische en cyclometrische functies.
\end{abstract}

\begin{document}

% FVD_CHAPTER_DATA_START
% Automatisch gegenereerd uit index.tex.
% Niet handmatig aanpassen.
\fvdchapterdata{3}{Afleiden van samengestelde functie}{2}
% FVD_CHAPTER_DATA_END






\maketitle


% ============================================================
% VOORKENNIS EN LEERDOELEN
% ============================================================

\begin{voorkennis}
  \item Je kunt veeltermfuncties afleiden.
  \item Je kunt de machtsregel gebruiken.
  \item Je kunt machten met rationale exponenten gebruiken.
  \item Je kunt de productregel en de quotiëntregel gebruiken.
  \item Je kunt algebraïsche uitdrukkingen vereenvoudigen.
  \item Je kunt een samengestelde functie herkennen als een functie
        die in een andere functie wordt ingevuld.
\end{voorkennis}

\begin{leerdoelen}
  \item Je kunt in een samengestelde functie de binnenste en buitenste
        functie herkennen.
  \item Je kunt een samengestelde functie schrijven als $f(g(x))$.
  \item Je kunt de kettingregel formuleren en gebruiken.
  \item Je kunt uitleggen waarom de afgeleide van $(f(g(x)))$
        niet alleen $f'(g(x))$ is.
  \item Je kunt samengestelde functies met meerdere lagen afleiden.
  \item Je kunt de kettingregel combineren met de productregel en
        de quotiëntregel.
  \item Je kunt bepalen welke afleidingsregel eerst moet worden toegepast.
  \item Je kunt een gevonden afgeleide algebraïsch vereenvoudigen.
  \item Je kunt domeinvoorwaarden blijven respecteren bij rationale
        en irrationale functies.
  \item Je kunt de structuur van een afgeleide analyseren en een
        geschikte afleidingsstrategie motiveren.
\end{leerdoelen}


% ============================================================
% 1. WAAROM EEN NIEUWE REGEL?
% ============================================================

\FVDsection{Wanneer functies in elkaar zitten}

Beschouw de functie
\[h(x)=(3x-1)^5.\]

De machtsregel
\[(x^n)'=nx^{n-1}\]
kunnen we hier niet rechtstreeks toepassen alsof er enkel $x$ tussen
de haakjes staat.

De functie bestaat namelijk uit twee opeenvolgende bewerkingen:

\begin{enumerate}
  \item eerst berekenen we
        \[u=3x-1;\]
  \item vervolgens berekenen we
        \[u^5.\]
\end{enumerate}

We kunnen dus schrijven:
\[g(x)=3x-1\]
en
\[f(u)=u^5.\]

Daarmee wordt
\[h(x)=f(g(x)).\]

De functie $g$ noemen we de \textbf{binnenste functie} en $f$ de
\textbf{buitenste functie}.


\begin{definitie}[Samengestelde functie]
Een functie van de vorm
\[h(x)=f(g(x))\]
noemen we een \textbf{samengestelde functie}.

We schrijven ook
\[h=f\circ g.\]

De functie $g$ wordt eerst toegepast en de functie $f$ daarna.
\end{definitie}


\begin{voorbeeld}[Binnen en buiten herkennen]
Beschouw
\[h(x)=\sqrt{x^2+4}.\]

We kunnen schrijven:
\[g(x)=x^2+4\]
en
\[f(u)=\sqrt{u}.\]

Dus:
\[h(x)=f(g(x)).\]

De binnenste functie is $x^2+4$ en de buitenste functie is de
wortelfunctie.
\end{voorbeeld}


% ============================================================
% 2. DE KETTINGREGEL
% ============================================================

\FVDsection{De kettingregel}

Wanneer
\[h(x)=f(g(x)),\]
dan verandert $h$ niet alleen omdat de buitenste functie $f$ verandert.

Ook het argument $g(x)$ verandert wanneer $x$ verandert.

Daarom moeten beide veranderingen in rekening worden gebracht.


\begin{definitie}[Kettingregel]
Als $g$ afleidbaar is in $x$ en $f$ afleidbaar is in $g(x)$, dan geldt:
\[\boxed{(f(g(x)))'=f'(g(x))\cdot g'(x).}\]
\end{definitie}


In woorden:

\begin{center}
\textbf{leid de buitenste functie af, laat de binnenkant staan,
en vermenigvuldig met de afgeleide van de binnenste functie.}
\end{center}


Een nuttig schema is:
\[\boxed{\text{buitenste afgeleide}\times\text{binnenste afgeleide}.}\]


% ============================================================
% 3. EERSTE VOORBEELD
% ============================================================

\FVDsection{Van buiten naar binnen}

Beschouw opnieuw
\[h(x)=(3x-1)^5.\]

We herkennen:

\begin{itemize}
  \item buitenste functie:
        \[f(u)=u^5;\]
  \item binnenste functie:
        \[g(x)=3x-1.\]
\end{itemize}

De afgeleiden zijn:
\[f'(u)=5u^4\]
en
\[g'(x)=3.\]

Met de kettingregel:
\[
h'(x)
=
5(3x-1)^4\cdot3.
\]

Dus:
\[\boxed{h'(x)=15(3x-1)^4.}\]


\begin{opmerking}
Een veelgemaakte fout is:
\[(3x-1)^5\longrightarrow5(3x-1)^4.\]

Daar ontbreekt de factor $3$.

De binnenste functie $3x-1$ verandert immers zelf ook met $x$:
\[(3x-1)'=3.\]

Daarom is:
\[\boxed{\bigl((3x-1)^5\bigr)'=5(3x-1)^4\cdot3.}\]
\end{opmerking}


% ============================================================
% 4. EEN PRAKTISCH STAPPENPLAN
% ============================================================

\FVDsection{Een praktisch stappenplan}

Bij een samengestelde functie kun je systematisch als volgt werken.

\begin{enumerate}
  \item Zoek de \textbf{buitenste bewerking}.
  \item Bepaal welke uitdrukking volledig binnen die bewerking zit.
  \item Leid de buitenste functie af terwijl je de binnenste uitdrukking
        voorlopig laat staan.
  \item Vermenigvuldig met de afgeleide van de binnenste functie.
  \item Herhaal de kettingregel wanneer de binnenste functie zelf opnieuw
        samengesteld is.
  \item Vereenvoudig het resultaat indien dat zinvol is.
\end{enumerate}


\begin{voorbeeld}[Wortelfunctie]
Bepaal de afgeleide van
\[f(x)=\sqrt{x^2+1}.\]

We schrijven:
\[f(x)=(x^2+1)^{1/2}.\]

De buitenste functie is een macht met exponent $\frac12$.
De binnenste functie is
\[u=x^2+1.\]

Dus:
\[
f'(x)
=
\frac12(x^2+1)^{-1/2}\cdot2x.
\]

Daaruit volgt:
\[\boxed{f'(x)=\frac{x}{\sqrt{x^2+1}}.}\]
\end{voorbeeld}


\begin{voorbeeld}[Negatieve exponent]
Bepaal de afgeleide van
\[f(x)=\frac{1}{(2x+3)^4}.\]

We herschrijven:
\[f(x)=(2x+3)^{-4}.\]

Met de kettingregel:
\[
f'(x)
=
-4(2x+3)^{-5}\cdot2.
\]

Dus:
\[\boxed{f'(x)=-\frac{8}{(2x+3)^5}.}\]
\end{voorbeeld}


% ============================================================
% 5. NOTATIE MET EEN TUSSENVARIABELE
% ============================================================

\FVDsection{Werken met een tussenvariabele}

Bij complexere functies kan een tussenvariabele de structuur duidelijker
maken.

Neem:
\[y=(x^2-3x+1)^6.\]

Stel:
\[u=x^2-3x+1.\]

Dan:
\[y=u^6.\]

We berekenen:
\[\frac{dy}{du}=6u^5\]
en
\[\frac{du}{dx}=2x-3.\]

Daarom:
\[
\frac{dy}{dx}
=
\frac{dy}{du}\cdot\frac{du}{dx}.
\]

Dus:
\[
\frac{dy}{dx}
=
6u^5(2x-3).
\]

We vervangen $u$ opnieuw:
\[\boxed{y'=6(x^2-3x+1)^5(2x-3).}\]


\begin{opmerking}
De notatie
\[\frac{dy}{dx}=\frac{dy}{du}\cdot\frac{du}{dx}\]
lijkt op het wegdelen van $du$.

Dat is een nuttige geheugensteun, maar de symbolen
$\frac{dy}{dx}$ en $\frac{dy}{du}$ zijn afgeleiden en geen gewone
breuken.

De eigenlijke wiskundige regel blijft:
\[\boxed{(f\circ g)'=(f'\circ g)\cdot g'.}\]
\end{opmerking}


% ============================================================
% 6. MEERDERE LAGEN
% ============================================================

\FVDsection{Samengestelde functies met meerdere lagen}

Een functie kan meer dan twee functielagen bevatten.

Beschouw:
\[f(x)=\sqrt{(2x-1)^3+4}.\]

We kunnen drie lagen onderscheiden:

\begin{enumerate}
  \item $2x-1$;
  \item $(2x-1)^3+4$;
  \item de vierkantswortel.
\end{enumerate}

We werken van buiten naar binnen.

Eerst:
\[
f'(x)
=
\frac{1}{2\sqrt{(2x-1)^3+4}}
\cdot
\left((2x-1)^3+4\right)'.
\]

Voor de resterende afgeleide gebruiken we opnieuw de kettingregel:
\[
\left((2x-1)^3+4\right)'
=
3(2x-1)^2\cdot2.
\]

Dus:
\[
f'(x)
=
\frac{1}{2\sqrt{(2x-1)^3+4}}
\cdot
6(2x-1)^2.
\]

Na vereenvoudigen:
\[\boxed{f'(x)=\frac{3(2x-1)^2}{\sqrt{(2x-1)^3+4}}.}\]


\begin{opmerking}
Bij meerdere lagen gebruiken we de kettingregel herhaaldelijk.

Je kunt dit onthouden als:
\[\boxed{\text{buiten}\longrightarrow\text{binnen}\longrightarrow
\text{nog verder naar binnen}.}\]

Van elke laag verschijnt uiteindelijk een afgeleide factor.
\end{opmerking}


% ============================================================
% 7. KETTINGREGEL OF PRODUCTREGEL?
% ============================================================

\FVDsection{Welke structuur zie je?}

Een belangrijke vaardigheid is herkennen welke regel bij welke structuur
hoort.

Vergelijk:
\[f(x)=(x^2+1)(x-3)\]
met
\[g(x)=(x^2+1)^3.\]

Bij $f$ worden twee functies met elkaar \textbf{vermenigvuldigd}.
Daarom gebruiken we de productregel.

Bij $g$ wordt de functie $x^2+1$ \textbf{ingevuld} in de machtsfunctie
$u^3$.
Daarom gebruiken we de kettingregel.


\begin{opmerking}
Haakjes betekenen niet automatisch dat je de kettingregel nodig hebt.

Vraag altijd:

\begin{center}
\textbf{Welke bewerking verbindt de functies?}
\end{center}

\begin{itemize}
  \item som of verschil $\longrightarrow$ som- of verschilregel;
  \item product $\longrightarrow$ productregel;
  \item quotiënt $\longrightarrow$ quotiëntregel;
  \item functie in een functie $\longrightarrow$ kettingregel.
\end{itemize}
\end{opmerking}


% ============================================================
% 8. REGELS COMBINEREN
% ============================================================

\FVDsection{De kettingregel combineren met andere regels}

In complexere functies moeten verschillende afleidingsregels gecombineerd
worden.

Beschouw:
\[f(x)=x^2(x+1)^4.\]

De buitenste structuur is een product:
\[x^2\cdot(x+1)^4.\]

We beginnen dus met de productregel:
\[
f'(x)
=
2x(x+1)^4
+
x^2\left((x+1)^4\right)'.
\]

Voor de tweede factor gebruiken we de kettingregel:
\[
\left((x+1)^4\right)'
=
4(x+1)^3.
\]

Dus:
\[
f'(x)
=
2x(x+1)^4
+
4x^2(x+1)^3.
\]

We kunnen ontbinden:
\[
f'(x)
=
2x(x+1)^3\bigl((x+1)+2x\bigr).
\]

Daarom:
\[\boxed{f'(x)=2x(x+1)^3(3x+1).}\]


\begin{voorbeeld}[Quotiëntregel en kettingregel]
Bepaal de afgeleide van
\[f(x)=\frac{(x^2+1)^3}{x-1}.\]

De buitenste structuur is een quotiënt.

Neem:
\[u(x)=(x^2+1)^3\]
en
\[v(x)=x-1.\]

Voor $u'$ gebruiken we de kettingregel:
\[u'(x)=3(x^2+1)^2\cdot2x=6x(x^2+1)^2.\]

Verder:
\[v'(x)=1.\]

Met de quotiëntregel:
\[
f'(x)
=
\frac{6x(x^2+1)^2(x-1)-(x^2+1)^3}{(x-1)^2}.
\]

We kunnen $(x^2+1)^2$ ontbinden:
\[
f'(x)
=
\frac{(x^2+1)^2\bigl(6x(x-1)-(x^2+1)\bigr)}
     {(x-1)^2}.
\]

Dus:
\[\boxed{f'(x)=
\frac{(x^2+1)^2(5x^2-6x-1)}
     {(x-1)^2}.}\]
\end{voorbeeld}


% ============================================================
% 9. DOMEIN BLIJFT BELANGRIJK
% ============================================================

\FVDsection{Afleiden verandert het oorspronkelijke domein niet}

Beschouw:
\[f(x)=\sqrt{5-2x}.\]

Voor de oorspronkelijke functie moet gelden:
\[5-2x\geq0.\]

Dus:
\[\operatorname{dom}f=\left]-\infty,\frac52\right].\]

Met de kettingregel:
\[
f'(x)
=
\frac{1}{2\sqrt{5-2x}}\cdot(-2).
\]

Dus:
\[\boxed{f'(x)=-\frac{1}{\sqrt{5-2x}}.}\]

Voor de afgeleide moet bovendien de noemer verschillend van nul zijn.
Daarom:
\[\operatorname{dom}f'=\left]-\infty,\frac52\right[.\]


\begin{opmerking}
Bij irrationale en rationale functies moet je steeds onderscheid maken
tussen:

\begin{itemize}
  \item het domein van de oorspronkelijke functie;
  \item het domein waarop de afgeleide bestaat.
\end{itemize}

Een correcte formule voor $f'$ is niet voldoende wanneer ze gebruikt
wordt buiten haar domein.
\end{opmerking}





% ============================================================
% 11. STRUCTUURBOOM
% ============================================================

\FVDsection{Een functie ontleden voor je afleidt}

Bij ingewikkeldere functies helpt het om eerst een structuurboom te maken.

Beschouw:
\[f(x)=\frac{\sqrt{x^2+1}}{(2x-3)^4}.\]

De hoofdstructuur is:
\[\boxed{\text{quotiënt}.}\]

De teller bevat:
\[\boxed{\text{wortel}\circ(x^2+1)}.\]

De noemer bevat:
\[\boxed{\text{vierde macht}\circ(2x-3)}.\]

Daaruit volgt de strategie:

\begin{enumerate}
  \item quotiëntregel;
  \item kettingregel voor de teller;
  \item kettingregel voor de noemer;
  \item algebraïsch vereenvoudigen.
\end{enumerate}


\begin{opmerking}
Bij complexe afgeleiden is de belangrijkste stap vaak niet het rekenen,
maar het herkennen van de structuur.

Een goede gewoonte is daarom:

\begin{center}
\textbf{eerst ontleden, dan afleiden.}
\end{center}
\end{opmerking}


% ============================================================
% 12. WETENSCHAPPELIJKE INTERPRETATIE
% ============================================================

\FVDsection{Een ketting van afhankelijkheden}

De kettingregel is niet alleen een rekentechniek.

In wetenschappelijke modellen hangt een grootheid vaak niet rechtstreeks
af van een andere grootheid, maar via een tussenliggende grootheid.

Stel bijvoorbeeld dat een grootheid $E$ afhangt van $r$:
\[E=f(r),\]
terwijl $r$ zelf verandert met de tijd:
\[r=g(t).\]

Dan is:
\[E(t)=f(g(t)).\]

De verandering van $E$ in de tijd wordt dan:
\[\boxed{\frac{dE}{dt}=\frac{dE}{dr}\cdot\frac{dr}{dt}.}\]

De kettingregel beschrijft dus hoe veranderingen zich door een
\textbf{keten van afhankelijkheden} voortplanten.


\begin{voorbeeld}[Oppervlakte van een groeiende cirkel]
De oppervlakte van een cirkel is
\[A=\pi r^2.\]

Stel dat de straal verandert volgens
\[r(t)=2t+1.\]

Dan:
\[A(t)=\pi(2t+1)^2.\]

Met de kettingregel:
\[
A'(t)
=
2\pi(2t+1)\cdot2.
\]

Dus:
\[\boxed{A'(t)=4\pi(2t+1).}\]

De afgeleide geeft aan hoe snel de oppervlakte verandert wanneer de
straal zelf met de tijd verandert.
\end{voorbeeld}




\begin{oefening}[Binnenste en buitenste functie]{\oefB\ESS}

Bepaal telkens de binnenste en buitenste functie.

\begin{enumerate}
  \item \[f(x)=(2x+3)^5.\]
  \item \[g(x)=\sqrt{x^2-4x+7}.\]
  \item \[h(x)=\frac{1}{(3x-1)^2}.\]
  \item \[p(x)=\left(x^3+2x\right)^4.\]
\end{enumerate}

\begin{oplossing}
\begin{enumerate}
  \item Binnen:
        \[u=2x+3.\]
        Buiten:
        \[f(u)=u^5.\]

  \item Binnen:
        \[u=x^2-4x+7.\]
        Buiten:
        \[g(u)=\sqrt{u}.\]

  \item Binnen:
        \[u=3x-1.\]
        Buiten:
        \[h(u)=u^{-2}.\]

  \item Binnen:
        \[u=x^3+2x.\]
        Buiten:
        \[p(u)=u^4.\]
\end{enumerate}
\end{oplossing}
\end{oefening}


\begin{oefening}[Kettingregel toepassen]{\oefB\ESS}

Bepaal de afgeleide.

\begin{enumerate}
  \item \[f(x)=(4x-1)^6.\]
  \item \[g(x)=(x^2+3)^4.\]
  \item \[h(x)=\sqrt{3x+5}.\]
  \item \[p(x)=\frac{1}{(x^2+1)^3}.\]
\end{enumerate}

\begin{oplossing}
\begin{enumerate}
  \item
  \[
  f'(x)=6(4x-1)^5\cdot4.
  \]
  Dus:
  \[\boxed{f'(x)=24(4x-1)^5.}\]

  \item
  \[
  g'(x)=4(x^2+3)^3\cdot2x.
  \]
  Dus:
  \[\boxed{g'(x)=8x(x^2+3)^3.}\]

  \item
  \[
  h'(x)=\frac{1}{2\sqrt{3x+5}}\cdot3.
  \]
  Dus:
  \[\boxed{h'(x)=\frac{3}{2\sqrt{3x+5}}.}\]

  \item We schrijven
  \[p(x)=(x^2+1)^{-3}.\]

  Dan:
  \[
  p'(x)
  =
  -3(x^2+1)^{-4}\cdot2x.
  \]

  Dus:
  \[\boxed{p'(x)=-\frac{6x}{(x^2+1)^4}.}\]
\end{enumerate}
\end{oplossing}
\end{oefening}


\begin{oefening}[Welke regel eerst?]{\oefU\ESS}

Bepaal telkens eerst de hoofdstructuur van de functie.
Geef aan welke afleidingsregel je eerst gebruikt en bereken daarna
de afgeleide.

\begin{enumerate}
  \item \[f(x)=x^2(x+1)^5.\]
  \item \[g(x)=\frac{(2x-1)^3}{x+2}.\]
  \item \[h(x)=\sqrt{x^2+4}.\]
  \item \[p(x)=\frac{1}{(x^2-1)^2}.\]
\end{enumerate}

\begin{oplossing}
\begin{enumerate}
  \item Hoofdstructuur: product.

  \[
  f'(x)
  =
  2x(x+1)^5
  +
  x^2\cdot5(x+1)^4.
  \]

  Dus:
  \[
  f'(x)
  =
  x(x+1)^4\bigl(2(x+1)+5x\bigr).
  \]

  \[\boxed{f'(x)=x(x+1)^4(7x+2).}\]

  \item Hoofdstructuur: quotiënt.

  De teller vereist de kettingregel:
  \[
  \bigl((2x-1)^3\bigr)'=6(2x-1)^2.
  \]

  Dus:
  \[
  g'(x)
  =
  \frac{6(2x-1)^2(x+2)-(2x-1)^3}{(x+2)^2}.
  \]

  Ontbinden geeft:
  \[
  g'(x)
  =
  \frac{(2x-1)^2\left(6(x+2)-(2x-1)\right)}
       {(x+2)^2}.
  \]

  Daarom:
  \[\boxed{g'(x)=\frac{(2x-1)^2(4x+13)}{(x+2)^2}.}\]

  \item Hoofdstructuur: samenstelling.

  \[
  h'(x)
  =
  \frac{1}{2\sqrt{x^2+4}}\cdot2x.
  \]

  Dus:
  \[\boxed{h'(x)=\frac{x}{\sqrt{x^2+4}}.}\]

  \item Hoofdstructuur: samenstelling na herschrijven.

  \[
  p(x)=(x^2-1)^{-2}.
  \]

  Dus:
  \[
  p'(x)
  =
  -2(x^2-1)^{-3}\cdot2x.
  \]

  Daarom:
  \[\boxed{p'(x)=-\frac{4x}{(x^2-1)^3}.}\]
\end{enumerate}
\end{oplossing}
\end{oefening}


\begin{oefening}[Een fout analyseren]{\oefU\ESS}

Een leerling bepaalt voor
\[f(x)=(x^2+1)^5\]
de afgeleide als
\[f'(x)=5(x^2+1)^4.\]

\begin{question}
Leg uit welke stap ontbreekt.
\end{question}

\begin{question}
Bereken de correcte afgeleide.
\end{question}

\begin{question}
Leg uit waarom deze fout vooral gevaarlijk wordt bij ingewikkeldere
samengestelde functies.
\end{question}

\begin{oplossing}
De leerling heeft alleen de buitenste macht afgeleid en is vergeten
dat de binnenste functie $x^2+1$ zelf ook verandert.

De binnenste afgeleide is:
\[(x^2+1)'=2x.\]

Daarom:
\[
f'(x)
=
5(x^2+1)^4\cdot2x.
\]

Dus:
\[\boxed{f'(x)=10x(x^2+1)^4.}\]

Bij functies met meerdere lagen moet voor elke functielaag de
bijbehorende afgeleide factor worden meegenomen. Wanneer één laag
wordt vergeten, is de volledige afgeleide fout.
\end{oplossing}
\end{oefening}



\end{document}